%% kapa: exam=1 round=26 year=2015 period=2 form=A subject=appraisal_law subject_ko=부동산관계법규 seq=3 number=3
%% source: Q-Net 감정평가사 기출문제 (cst00309) artlSeq=5138001 file=extracted/1008/5138001/제26회_감정평가사 1차_기출/제26회_감정평가사 1차_기출/2교시_감평_A형.pdf
%% conversion: kordoc 4.24.51 -> tools/gen1.py (auto); review=proofread-sonnet; flags=none
\begin{question}{3}
국토의 계획 및 이용에 관한 법령상 입지규제최소구역으로 지정할 수 있는 곳을 모두 고른 것은?
\begin{framedtext}
ㄱ. 도시ㆍ군기본계획에 따른 생활권의 중심지역 \par
ㄴ. 세 개 이상의 노선이 교차하는 대중교통 결절지로부터 1킬로미터 이내에 위치한 지역 \par
ㄷ.「도시 및 주거환경정비법」에 따른 노후ㆍ불량건축물이 밀집한 주거지역으로 정비가 시급한 지역 \par
ㄹ.「도시재생 활성화 및 지원에 관한 특별법」에 따른 도시재생활성화지역 중 도시경제기반형 활성화계획을 수립하는 지역 \par
\end{framedtext}
\begin{choices}
  \item ㄱ, ㄹ
  \item ㄴ, ㄷ
  \item ㄱ, ㄴ, ㄷ
  \item ㄴ, ㄷ, ㄹ
  \item ㄱ, ㄴ, ㄷ, ㄹ
\end{choices}
\end{question}
\answer{5}
